\documentclass[aspectratio=169,11pt]{beamer}

% --- minimalni, cist stil ---
\usetheme{default}
\usecolortheme{dove}
\usefonttheme{professionalfonts}
\setbeamertemplate{navigation symbols}{}
\setbeamertemplate{footline}[frame number]
\setbeamertemplate{itemize items}[circle]
\setbeamertemplate{blocks}[rounded][shadow=false]

\usepackage[T1]{fontenc}
\usepackage[utf8]{inputenc}
\usepackage{amsmath,amssymb}
\usepackage{booktabs}
\usepackage{listings}
\usepackage{tikz}
\usetikzlibrary{arrows.meta,positioning}

% oznake klasa detektora
\newcommand{\Dp}{\ensuremath{\mathcal{P}}}
\newcommand{\Ds}{\ensuremath{\mathcal{S}}}
\newcommand{\Dw}{\ensuremath{\mathcal{W}}}
\newcommand{\Dv}{\ensuremath{\mathcal{V}}}
\newcommand{\eP}{\ensuremath{\Diamond\mathcal{P}}}
\newcommand{\eS}{\ensuremath{\Diamond\mathcal{S}}}
\newcommand{\eW}{\ensuremath{\Diamond\mathcal{W}}}
\newcommand{\eV}{\ensuremath{\Diamond\mathcal{V}}}
\newcommand{\Crashed}{\mathrm{Crashed}}
\newcommand{\Up}{\mathrm{Up}}

\lstset{
  basicstyle=\ttfamily\scriptsize,
  keywordstyle=\bfseries,
  commentstyle=\itshape\color{gray!60!black},
  frame=single,breaklines=true,
  showstringspaces=false,
  language=Java,
  literate={č}{{\v{c}}}1 {ć}{{\'c}}1 {ž}{{\v{z}}}1 {š}{{\v{s}}}1 {đ}{{\dj}}1
           {Č}{{\v{C}}}1 {Ć}{{\'C}}1 {Ž}{{\v{Z}}}1 {Š}{{\v{S}}}1 {Đ}{{\DJ}}1
}

\title{Otkrivači grešaka u distribuiranim sustavima}
\author{Bartol Gracin \and Vid Joha \and Matija Njegač}
\institute{Distribuirani procesi}
\date{}

\begin{document}

\begin{frame}
  \titlepage
\end{frame}

\begin{frame}{Pregled}
  \begin{enumerate}
    \item Motivacija: asinkroni sustavi i nemogućnost konsenzusa
    \item Otkrivači grešaka: potpunost i točnost
    \item Osam klasa i reducibilnost
    \item Konsenzus pomoću otkrivača (\Ds{} i \eS)
    \item Implementacije:
      \begin{itemize}
        \item Algoritam 15.3 --- konsenzus uz \Ds
        \item Algoritam 15.4 --- konsenzus uz \eS
        \item Algoritam 15.7 --- otkrivač \eP
      \end{itemize}
  \end{enumerate}
\end{frame}

% =======================  TEORIJA  =======================

\section{Teorija}

\begin{frame}{Motivacija}
  \begin{itemize}
    \item \textbf{Asinkroni model:} nema gornje granice na trajanje koraka ni na
      kašnjenje poruke.
    \item Proces može zakazati \emph{zaustavljanjem} (crash) --- trajno stane.
    \item Problem: ne razlikujemo \emph{srušen} proces od \emph{sporog}.
    \item \textbf{FLP:} u asinkronom sustavu konsenzus je nemoguć već uz jedan kvar.
  \end{itemize}
  \vfill
  \begin{block}{Rješenje}
    Sustavu dodajemo \emph{otkrivač grešaka} --- orakl koji daje (nepouzdane)
    naznake o tome tko se srušio. Time konsenzus postaje rješiv.
  \end{block}
\end{frame}

\begin{frame}{Otkrivači grešaka}
  \begin{itemize}
    \item Svaki proces $p_i$ ima lokalni modul $D_i$ koji prati ostale procese.
    \item $D_p(t)$ = skup procesa za koje $p$ u trenutku $t$ sumnja da su se
      srušili.
    \item Sumnja se temelji na tome koliko dugo nismo čuli neki proces.
    \item \textbf{Nepouzdani su:} ispravan proces može biti osumnjičen, pa kasnije
      maknut s popisa.
  \end{itemize}
  \vfill
  Kvalitetu opisujemo dvama svojstvima: \textbf{potpunost} i \textbf{točnost}.
\end{frame}

\begin{frame}{Potpunost}
  Kada otkrivač \emph{mora} posumnjati na srušeni proces.
  \begin{block}{Jaka potpunost}
    Nakon nekog trenutka svaki srušeni proces trajno je sumnjiv \emph{svakom}
    ispravnom procesu.
  \end{block}
  \begin{block}{Slaba potpunost}
    Nakon nekog trenutka svaki srušeni proces trajno je sumnjiv \emph{nekom}
    ispravnom procesu.
  \end{block}
  \vfill
  \footnotesize Sama potpunost nije dovoljna (detektor koji sumnja na sve
  zadovoljava jaku potpunost, a beskoristan je).
\end{frame}

\begin{frame}{Točnost}
  Kada otkrivač \emph{smije} pogriješiti (posumnjati na ispravan proces).
  \begin{itemize}
    \item \textbf{Jaka:} nijedan ispravan proces nikada nije sumnjiv.
    \item \textbf{Slaba:} neki ispravan proces nikada nije sumnjiv.
    \item \textbf{Konačna jaka:} postoji trenutak nakon kojeg nijedan ispravan
      proces nije sumnjiv.
    \item \textbf{Konačna slaba:} postoji trenutak nakon kojeg neki ispravan
      proces nije sumnjiv.
  \end{itemize}
  \vfill
  \footnotesize ,,Konačna'' varijanta dopušta rane greške, ali zahtijeva da one
  \emph{prestanu}.
\end{frame}

\begin{frame}{Osam klasa otkrivača}
  \centering
  \small
  \begin{tabular}{lllc}
    \toprule
    Klasa & Potpunost & Točnost & Oznaka \\
    \midrule
    Savršeni             & jaka  & jaka          & \Dp \\
    Konačno savršeni     & jaka  & konačna jaka  & \eP \\
    Jaki                 & jaka  & slaba         & \Ds \\
    Konačno jaki         & jaka  & konačna slaba & \eS \\
    Slabi                & slaba & slaba         & \Dw \\
    Konačno slabi        & slaba & konačna slaba & \eW \\
    Slaba potpunost i jaka točnost                  & slaba & jaka          & \Dv \\
    Slaba potpunost i konačna jaka točnost                  & slaba & konačna jaka  & \eV \\
    \bottomrule
  \end{tabular}
\end{frame}

\begin{frame}{Reducibilnost}
  Transformacijom slabe potpunosti u jaku dobivamo ekvivalencije:
  \[
    \Dp \equiv \Dv,\quad \Ds \equiv \Dw,\quad \eP \equiv \eV,\quad \eS \equiv \eW.
  \]
  \vfill
  \centering
  \begin{tikzpicture}[>=Latex, node distance=10mm and 20mm, every node/.style={font=\small}]
    \node (P) {$\Dp \equiv \Dv$};
    \node (eP) [below left=of P] {$\eP \equiv \eV$};
    \node (S) [below right=of P] {$\Ds \equiv \Dw$};
    \node (eS) [below right=of eP] {$\eS \equiv \eW$};
    \draw[->] (P) -- (eP); \draw[->] (P) -- (S);
    \draw[->] (eP) -- (eS); \draw[->] (S) -- (eS);
  \end{tikzpicture}
  \vfill
  \footnotesize Za konsenzus je dovoljno riješiti dvije klase: \Ds{} i \eS.
\end{frame}

\begin{frame}{Konsenzus pomoću otkrivača}
  Svaki proces predlaže vrijednost; traži se \emph{konačnost}, \emph{usuglašenost}
  i \emph{valjanost}.
  \vfill
  \begin{block}{Uz \Ds{} (jaki)}
    Konsenzus rješiv \textbf{bez ograničenja na broj kvarova} (do $n-1$). \\
    $\Rightarrow$ Algoritam 15.3.
  \end{block}
  \begin{block}{Uz \eS{} (konačno jaki)}
    Konsenzus rješiv \textbf{samo ako je većina ispravna} ($f < n/2$). \\
    $\Rightarrow$ Algoritam 15.4.
  \end{block}
  \vfill
  \footnotesize \eS{} je najslabiji otkrivač koji uopće rješava konsenzus; u praksi
  gradimo barem \eP{} (Algoritam 15.7).
\end{frame}

% =======================  ALGORITAM 15.3  =======================

\section{Algoritam 15.3}

\begin{frame}{Algoritam 15.3 --- konsenzus uz \Ds}
  \textbf{Ideja.} Konsenzus koji podnosi \textbf{bilo koji broj kvarova} (do
  $n-1$) i ne treba većinu; radi u tri faze.
  \begin{itemize}
    \item Svaki proces drži vektor $V_p$ --- svoje procjene vrijednosti koje su
      predložili svi procesi (na početku prazne, osim vlastite).
    \item $\Delta_p$ sadrži vrijednosti naučene u zadnjoj rundi.
    \item Upit ,,$q \in D_p$'' (koga sumnjičimo) ostvaruje ugrađeni heartbeat
      detektor, poruke ,,Alive''.
  \end{itemize}
\end{frame}

\begin{frame}{Algoritam 15.3 --- opis}
  \begin{itemize}
    \item \textbf{Faza 1} ($n-1$ rundi). U svakoj rundi proces svima pošalje
      vrijednosti koje je zadnje naučio i čeka poruku te runde od svakog procesa
      kojeg \emph{ne} sumnjiči --- osumnjičene preskače. Svaku novu vrijednost
      upiše u $V_p$. Nakon $n-1$ rundi svi ispravni imaju isti skup
      ,,preživjelih'' vrijednosti.
    \item \textbf{Faza 2.} Procesi razmijene cijele vektore. Ako je bilo koji
      primljeni vektor imao prazno mjesto, i vlastito se mjesto isprazni --- tako
      ispadaju vrijednosti oko kojih nema potpune sloge. Barem jedno mjesto
      ostaje popunjeno.
    \item \textbf{Faza 3.} Svi odluče \emph{prvu} popunjenu vrijednost vektora;
      ona je ista kod svih.
  \end{itemize}
\end{frame}

\begin{frame}[fragile,allowframebreaks]{Algoritam 15.3 --- kod}
\begin{lstlisting}
for(int runda=1; runda < numProc; runda++){
    // Faza 1: posalji svoju Deltu svima
    for (int j = 0; j < numProc; j++)
        if (j != myId) process.sendMsg(j,"Consensus",runda,process.getDelta());

    // cekaj dok ne primis od svih ili ih ne osumnjicis (upit detektoru)
    do {
        for (int j = 0; j < numProc; j++)
            if (j != myId){ process.sendMsg(j,"Alive",runda,process.getDelta());
                            process.provjeriProces(j); }
        Thread.sleep(1000);
    } while(process.kolikoSamPorukaPrimioURundi(runda)
            + process.kolikoJeProcesaOsumnjiceno() < numProc-1);

    process.obradiPoruke(runda);   // spoji primljene vrijednosti u V
    process.novaRunda();
}
\end{lstlisting}

\begin{lstlisting}
// Faza 1: za svako mjesto i koje je jos null, uzmi vrijednost iz poruke
public void obradiPoruke(int runda){
    for (int i = 0; i < numProc; i++) Delta.set(i, null);
    for (int i = 0; i < numProc; i++){
        if(V.get(i)==null){
            for(int j = 0; j < Poruke.get(runda-1).size(); j++){
                String[] t = Poruke.get(runda-1).get(j).getMessage().trim().split("\\s+");
                Integer val = t[i+1].equals("null") ? null : Integer.parseInt(t[i+1]);
                if(val!=null){ V.set(i,val); Delta.set(i,val); }
            }
        }
    }
}

// Faza 2: ako je neki primljeni V_q imao null na mjestu i, i mi ga ponistimo
public void provjeriZadnjePoruke(int proces){
    for(int i = 0; i<ZadnjePoruke.size(); i++){
        if(ZadnjePoruke.get(i).getSrcId()==proces){
            String[] t = ZadnjePoruke.get(i).getMessage().trim().split("\\s+");
            Integer val = t[i+1].equals("null") ? null : Integer.parseInt(t[i+1]);
            if(val==null) V.set(i, null);
        }
    }
}
\end{lstlisting}
\end{frame}

\begin{frame}{Algoritam 15.3 --- zašto radi}
  \begin{itemize}
    \item \textbf{Otpornost:} nijedna faza ne traži većinu --- čeka se svaki
      neosumnjičeni proces, osumnjičeni se preskaču.
    \item \textbf{Usuglašenost:} slaba točnost \Ds{} jamči ispravan proces
      $p^\ast$ kojeg nitko ne sumnjiči; njegova vrijednost preživi u svim
      vektorima $\Rightarrow$ svi odluče isto.
    \item \textbf{Konačnost:} jaka potpunost $\Rightarrow$ srušeni se na kraju
      osumnjiče, pa čekanje ne blokira zauvijek.
  \end{itemize}
\end{frame}

% =======================  ALGORITAM 15.4  =======================

\section{Algoritam 15.4}

\begin{frame}{Algoritam 15.4 --- konsenzus uz \eS}
  \textbf{Ideja.} Konsenzus \emph{rotirajućim koordinatorom}: u rundi $r$
  koordinator je $c=(r \bmod n)+1$, i sve poruke idu prema koordinatoru ili od
  njega.
  \begin{itemize}
    \item Traži \textbf{ispravnu većinu} procesa ($f < n/2$).
    \item Svaki proces pamti procjenu odluke $estimate_p$ i vremensku oznaku
      $ts_p$ (rundu u kojoj je procjenu zadnji put promijenio).
    \item Sumnju na koordinatora daje otkrivač \eS.
  \end{itemize}
\end{frame}

\begin{frame}{Algoritam 15.4 --- opis}
  U svakoj rundi koordinator pokušava ,,zaključati'' zajedničku vrijednost kroz
  četiri faze:
  \begin{itemize}
    \item \textbf{Faza 1.} Svi procesi pošalju koordinatoru svoju procjenu i
      oznaku $(estimate_p, ts_p)$.
    \item \textbf{Faza 2.} Koordinator pričeka procjene od većine, odabere
      najsvježiju (najveća oznaka $ts$) i pošalje je svima kao prijedlog.
    \item \textbf{Faza 3.} Tko primi prijedlog, prihvati ga i odgovori
      \emph{ack}; tko posumnja na koordinatora, odgovori \emph{nack}.
    \item \textbf{Faza 4.} Dobije li koordinator \emph{ack} od većine, vrijednost
      je zaključana i emitira \emph{decide}; svi koji ga prime odluče tu
      vrijednost.
  \end{itemize}
  \vfill
  \footnotesize Sruši li se koordinator ili je osumnjičen, kreće iduća runda s
  novim koordinatorom.
\end{frame}

\begin{frame}[fragile,allowframebreaks]{Algoritam 15.4 --- kod}
\begin{lstlisting}
public void runConsensus(){
    while(state == 0){
        receivedPhase1Messages.clear();
        receivedPhase2Messages.clear();
        round++;
        int coordinatorId = (round%n);
        phase1_sendEstimate(coordinatorId);
        phase2_coordinatorSelect(coordinatorId);
        phase3_ackNack(coordinatorId);
        phase4_decide(coordinatorId);
        try { Thread.sleep(1000); } catch (InterruptedException e) {}
    }
    System.out.println("FINAL DECISION = " + estimate);
}
\end{lstlisting}

\begin{lstlisting}
public void phase1_sendEstimate(int coordinatorId) {
    String payload = estimate + " " + estimateRound;   // estimate + timestamp
    if (coordinatorId != myId)
        linker.sendMsg(coordinatorId, "PHASE1", payload);
}

public void phase2_coordinatorSelect(int coordinatorId) {
    if (myId != coordinatorId) return;          // samo koordinator
    int bestEstimate = estimate, bestRound = estimateRound;
    int majority = (n/2) + 1;
    while(receivedPhase1Messages.size() < majority-1)
        try { Thread.sleep(100); } catch (InterruptedException e){ break; }
    for (Msg m : new ArrayList<>(receivedPhase1Messages)) {
        String[] parts = m.getMessage().trim().split("\\s+");
        int value = Integer.parseInt(parts[0]), ts = Integer.parseInt(parts[1]);
        if (ts > bestRound) { bestRound = ts; bestEstimate = value; }  // najsvjeziji
    }
    estimate = bestEstimate; estimateRound = round;
    linker.broadcastToAll("PHASE2", estimate + " " + estimateRound);
    receivedPhase2Messages.add(new Msg(myId, myId, "PHASE2", estimate+" "+estimateRound));
}
\end{lstlisting}

\begin{lstlisting}
public void phase3_ackNack(int coordinatorId) {
    Msg coordinatorMsg = null;
    while(receivedPhase2Messages.isEmpty())
        try { Thread.sleep(100); } catch(InterruptedException e){}
    for (Msg m : new ArrayList<>(receivedPhase2Messages)) {
        if (m.getSrcId()==coordinatorId && m.getDestId()==myId) {
            coordinatorMsg = m;
            String[] parts = m.getMessage().trim().split("\\s+");
            estimate = Integer.parseInt(parts[0]); estimateRound = round;  // prihvati
            break;
        }
    }
    if (coordinatorMsg != null || coordinatorId == myId) {   // ACK
        if (coordinatorId != myId) linker.sendMsg(coordinatorId, "ACK", round + "");
        else receivedPhase3Replies.add(new Msg(myId, myId, "ACK", round + ""));
    } else {                                                 // NACK
        if (coordinatorId != myId) linker.sendMsg(coordinatorId, "NACK", round + "");
    }
}

public void phase4_decide(int coordinatorId) {
    if (myId != coordinatorId) return;         // samo koordinator odlucuje
    int ackCount = 0, majority = (n/2) + 1;
    while(receivedPhase3Replies.size() < majority)
        try { Thread.sleep(100); } catch(InterruptedException e){}
    for (Msg m : new ArrayList<>(receivedPhase3Replies))
        if (m.getSrcId()!=coordinatorId && m.getTag().equals("ACK")) ackCount++;
    if (ackCount >= majority) {                 // vecina prihvatila -> DECIDE
        linker.broadcastToAll("DECIDE", estimate + " " + round);
        state = 1;
    }
    receivedPhase3Replies.clear();
}
\end{lstlisting}

\begin{lstlisting}
// obrada odluke u handleMsg
else if (tag.equals("DECIDE")) {
    String[] parts = m.getMessage().trim().split("\\s+");
    estimate = Integer.parseInt(parts[0]);
    state = 1;
}
\end{lstlisting}
\end{frame}

\begin{frame}{Algoritam 15.4 --- zašto radi}
  \begin{itemize}
    \item \textbf{Zašto većina:} faze 2 i 4 čekaju $\lceil(n{+}1)/2\rceil$
      odgovora; dvije se većine uvijek preklapaju $\Rightarrow$ zaključana
      vrijednost prenosi se u sve kasnije runde (usuglašenost).
    \item Bez ispravne većine čekanje bi moglo blokirati $\Rightarrow$ uvjet
      $f<n/2$.
    \item \textbf{Konačnost:} konačna slaba točnost \eS{} $\Rightarrow$ na kraju
      postoji ispravan, neosumnjičen koordinator koji zaključa vrijednost i
      emitira \textit{decide}.
  \end{itemize}
\end{frame}

% =======================  ALGORITAM 15.7  =======================

\section{Algoritam 15.7}

\begin{frame}{Algoritam 15.7 --- otkrivač \eP}
  \textbf{Ideja.} Savršen otkrivač \Dp{} nije ostvariv u asinkronom sustavu;
  najbolje što možemo jest da greške \emph{na kraju} prestanu --- to je \eP.
  \begin{itemize}
    \item Radi u \emph{parcijalnoj sinkroniji} (nakon nepoznatog trenutka sustav
      postane sinkron).
    \item Vrijeme se mjeri lokalnim satom (otkucaji), ne stvarnim satom.
    \item Ključ: \textbf{adaptivni timeout} $\Delta_p(q)$ koji raste nakon svake
      pogrešne sumnje.
  \end{itemize}
\end{frame}

\begin{frame}{Algoritam 15.7 --- opis (tri zadatka)}
  Tri paralelna zadatka; vrijeme mjeri lokalni sat koji broji otkucaje.
  \begin{itemize}
    \item \textbf{Zadatak 1 (heartbeat).} Proces periodički svima šalje poruku
      ,,živ sam''.
    \item \textbf{Zadatak 2 (timeout).} Ako od procesa $q$ nije čuo ništa dulje od
      $\Delta_p(q)$ otkucaja, doda ga među sumnjive.
    \item \textbf{Zadatak 3 (kajanje).} Primi li poruku od procesa kojeg
      sumnjiči, makne ga s popisa i \emph{poveća} $\Delta_p(q)$ za jedan.
  \end{itemize}
  \vfill
  \footnotesize Timeout se sam prilagođava: raste dok ne prijeđe stvarno
  kašnjenje, nakon čega lažne sumnje prestaju.
\end{frame}

\begin{frame}[fragile,allowframebreaks]{Algoritam 15.7 --- kod}
\begin{lstlisting}
// Zadatak 1: posalji "ziv sam" svima
public void sendHeartbeats() {
    if (crashed) return;
    broadcastMsg("alive", (int) clock);
}

// Zadatak 2: istek vremena -> sumnja
public synchronized void checkTimeouts() {
    for (int q = 0; q < N; q++) {
        if (q == myId) continue;
        if (!Output.contains(q) && (clock - lastAlive[q]) > delta[q]) {
            Output.add(q);
        }
    }
}

// Zadatak 3: kad primim "ziv sam" od q -> kajanje + veci timeout
public synchronized void handleMsg(Msg m, int q, String tag) {
    if (tag.equals("alive")) {
        lastAlive[q] = clock;
        if (Output.contains(q)) {
            Output.remove(q);
            delta[q] = delta[q] + 1;
        }
    }
}
\end{lstlisting}
\end{frame}

\begin{frame}{Algoritam 15.7 --- ispravnost i test}
  \textbf{Jaka potpunost.} Srušeni proces prestaje slati; tišina raste iznad
  $\Delta_p(q)$ $\Rightarrow$ trajno osumnjičen.

  \textbf{Konačna jaka točnost.} Nakon sinkronosti $\Delta$ prerasta stvarno
  kašnjenje $\Rightarrow$ pogrešne sumnje prestaju.
  \vfill
  \centering
  \small
  \begin{tabular}{ll}
    \toprule
    Scenarij & Rezultat \\
    \midrule
    kvar $p_2$ (t=10)        & svi trajno sumnjaju na $p_2$; potpunost \\
    spor $p_2$ (svakih 8)    & $\Delta$: $5{\to}6{\to}7{\to}8$, sumnje prestanu; točnost \\
    \bottomrule
  \end{tabular}
\end{frame}

% =======================  ZAKLJUCAK  =======================

\begin{frame}{Zaključak}
  \begin{itemize}
    \item Otkrivači se opisuju \textbf{potpunošću} i \textbf{točnošću} $\Rightarrow$
      osam klasa.
    \item Zbog reducibilnosti za konsenzus su bitne \Ds{} i \eS.
    \item \textbf{15.3} (\Ds): konsenzus uz bilo koji broj kvarova.
    \item \textbf{15.4} (\eS): konsenzus uz ispravnu većinu (rotirajući koordinator).
    \item \textbf{15.7} (\eP): ostvariv otkrivač u parcijalnoj sinkroniji
      (adaptivni timeout).
  \end{itemize}
  \vfill
  \centering \emph{Hvala na pažnji!}
\end{frame}

\end{document}
